\subsection{PROPERTIES OF QUANTIFIERS}

From now on we shall have to consider only quantified theories. For the rest of this section, \(\mathfrak{C}\) will denote such a theory, and \(\mathfrak{C}_0\) will denote the theory without explicit axioms which has the same signs as \(\mathfrak{C}\) and whose only schemes are S1 through S5. \(\mathfrak{C}\) is stronger than \(\mathfrak{C}_0\) .

C29. Let R be a relation in C and let x be a letter. Then the relations ``not \((\exists x)R\) '' and \((\forall x)\) (not R) are equivalent in C.

It is sufficient to prove the criterion in the theory \(\mathfrak{G}_0\) in which \(x\) is not a constant. The theorem \(R \Longleftrightarrow\) (not not \(R\) ) gives us, by C3 ( \(\S 2\) , no. 3), the theorems

\[
(\exists x) R \Longrightarrow (\tau_ {x} (R) | x) (\text { not   not } R)
\]

and

\[
(\exists x) (\text { not   not } R) \Longrightarrow (\tau_ {x} (\text { not   not } R) | x) R.
\]

Applying S5, we deduce the theorems in \(\mathfrak{C}_{0}\)

\[
(\exists x) R \Longrightarrow (\exists x) (\text { not   not } R), \quad (\exists x) (\text { not   not } R) \Longrightarrow (\exists x) R,
\]

whence the theorem \((\exists x)R \Longleftrightarrow (\exists x)\) (not not \(R\) ). Now

\[
(\exists x) (\text { not   not } R)
\]

is equivalent in \(\mathcal{G}_0\) to ``not not \((\exists x)\) (not not \(R\) )'', that is to ``not \((\forall x)\) (not \(R\) )''. Hence the result.

The criteria C28 and C29 enable us to deduce properties of one of the quantifiers from properties of the other.

C30. Let R be a relation in G, let T be a term in G, and let x be a letter. Then the relation \((\forall x)R \Longrightarrow (T|x)R\) is a theorem in G.

By S5, \((T|x)\) (not \(R\) ) \(\Longrightarrow\) ( \(\tau_{x}\) (not \(R\) ) \(|x\) ) (not \(R\) ) is an axiom. This relation is identical with

\[
(\text { not } (T | x) R) \Longrightarrow \text { not } ((\tau_ {x} (\text { not } R) | x) R).
\]

Hence \((\tau_{\mathbf{x}}(\text{not } R)|x)R \Longrightarrow (T|x)R\) is a theorem in \(\mathfrak{C}\) . Now use C26 (no. 1).

Let \(\pmb{R}\) be a relation in \(\mathfrak{C}\) . By C26, C27, and C30 it is the same (provided the letter \(\pmb{x}\) is not a constant of \(\mathfrak{C}\) ) whether we state the theorem \(\pmb{R}\) in \(\mathfrak{C}\) , or the theorem \((\forall x)\pmb{R}\) , or the metamathematical rule: if \(\pmb{T}\) is any term in \(\mathfrak{C}\) , then \((\pmb{T}|\pmb{x})\pmb{R}\) is a theorem in \(\mathfrak{C}\) .

C31. Let R and S be relations in C, and let x be a letter which is not a constant of C. If \(R \Longrightarrow S\) (resp. \(R \Longleftrightarrow S\) ) is a theorem in C, then

\[
\begin{array}{c c} (\forall x) R \Longrightarrow (\forall x) S, & (\exists x) R \Longrightarrow (\exists x) S \\ \text{[resp. } (\forall x) R \Longleftrightarrow (\forall x) S, & (\exists x) R \Longleftrightarrow (\exists x) S \text{]} \end{array}
\]

are theorems in \(\mathfrak{G}\) .

Suppose that \(R \Longrightarrow S\) is a theorem in \(\mathcal{G}\) . Let us adjoin the hypothesis \((\forall x)R\) (in which \(x\) does not appear). Then \(R\) , hence \(S\) , and therefore also \((\forall x)S\) , are true. Consequently \((\forall x)R \Longrightarrow (\forall x)S\) is a theorem in \(\mathcal{G}\) . It follows that if \(R \Longleftrightarrow S\) is a theorem in \(\mathcal{G}\) , then so is

\[
(\forall x) R \Longleftrightarrow (\forall x) S.
\]

The rules relating to \(\exists\) can now be deduced by using C29.

C32. Let R and S be relations in C, and let x be a letter. Then the relations

\[
\begin{array}{l} (\forall x) (R \text {   and   } S) \iff ((\forall x) R \text {   and   } (\forall x) S), \\ (\exists x) (R \text {   or   } S) \iff ((\exists x) R \text {   or   } (\exists x) S) \end{array}
\]

are theorems in \(\mathfrak{C}\) .

It is sufficient to prove these criteria in \(C_{0}\) , in which x is not a constant. If \((\forall x)(R \text{ and } S)\) is true, then ``R and \(S\)'' is true, and therefore each of the relations R, S is true. Consequently each of the relations \((\forall x)R\) , \((\forall x)S\) is true, and hence `` \((\forall x)R\) and \((\forall x)S\)'' is true. Similarly one shows that if `` \((\forall x)R\) and \((\forall x)S\)'' is true, then \((\forall x)(R \text{ and } S)\) is true. Hence the first theorem. The second follows by applying C29.

It should be noted that if \((\forall x)(R \text{ or } S)\) is a theorem in G, we may not conclude that \(((\forall x)R \text{ or } (\forall x)S)\) is a theorem in G. Intuitively, to say that the relation \((\forall x)(R \text{ or } S)\) is true means that for each object x, at least one of the relations R, S is true; but in general only one of the two will be true, and whether it is R or S will depend on the choice of x. Likewise, if \(((\exists x)R \text{ and } (\exists x)S)\) is a theorem in G, we may not conclude that \((\exists x)(R \text{ and } S)\) is a theorem in G. However, there is the following criterion:

C33. Let R and S be relations in G, and let x be a letter which does not appear in R. Then the relations

\[
\begin{array}{l l} (\forall x) (R \text {   or   } S) & \Longleftrightarrow (R \text {   or   } (\forall x) S), \\ (\exists x) (R \text {   and   } S) & \Longleftrightarrow (R \text {   and   } (\exists x) S) \end{array}
\]

are theorems in \(\mathfrak{C}\) .

It is sufficient to establish the criterion in \(\mathfrak{C}_0\) , in which \(x\) is not a constant. Let \(\mathfrak{C}'\) be the theory obtained by adjoining \((\forall x)(R\) or \(S)\) to the axioms of \(\mathfrak{C}_0\) . In \(\mathfrak{C}'\) , `` \(R\) or \(S\) '', and therefore (not \(R) \Longrightarrow S\) , are theorems. If ``not \(R\) '' is true (a hypothesis in which \(x\) does not feature), then \(S\) and therefore also \((\forall x)S\) are true. Consequently

\[
(\text { not } R) \Longrightarrow (\forall x) S
\]

is a theorem in \(\mathcal{C}'\) , and hence \((\forall x)(R\) or \(S) \Longrightarrow (R\) or \((\forall x)S)\) is a theorem in \(\mathcal{C}_0\) . Likewise, if `` \(R\) or \((\forall x)S\) '' is true, then `` \(R\) or \(S\) '' and therefore \((\forall x)(R\) or \(S)\) are true. Consequently

\[
(R \text {   or   } (\forall x) S) \Longrightarrow (\forall x) (R \text {   or   } S)
\]

is a theorem in \(\mathfrak{G}_0\) . The rule relating to \(\exists\) follows by applying C29.

C34. Let R be a relation and let x and y be letters. Then the relations

\[
\begin{array}{l} (\forall x) (\forall y) R \iff (\forall y) (\forall x) R, \\ (\exists x) (\exists y) R \iff (\exists y) (\exists x) R, \\ (\exists x) (\forall y) R \implies (\forall y) (\exists x) R \end{array}
\]

are theorems in \(\mathfrak{G}\) .

It is sufficient to prove these theorems in \(\mathcal{G}_0\) , in which \(x\) and \(y\) are not constants. If \((\forall x)(\forall y)R\) is true, then \((\forall y)R\) , and therefore \(R\) , hence \((\forall x)R\) , hence \((\forall y)(\forall x)R\) , are true. Likewise, if \((\forall y)(\forall x)R\) is true, then \((\forall x)(\forall y)R\) is true; and the first theorem follows. The second now follows by use of C29. Finally, since \((\forall y)R \Longrightarrow R\) is a theorem in \(\mathcal{C}_0\) , so is \((\exists x)(\forall y)R \Longrightarrow (\exists x)R\) by C31; if \((\exists x)(\forall y)R\) is true, then \((\exists x)R\) is true, and therefore so is \((\forall y)(\exists x)R\) . Hence the third theorem.

On the other hand, if \((\forall y)(\exists x)R\) is a theorem in C, we may not conclude that \((\exists x)(\forall y)R\) is a theorem in C. Intuitively, to say that the relation \((\forall y)(\exists x)R\) is true means that, given any object y, there is an object x such that R is a true relation between the objects x and y. But in general the object x will depend on the choice of the object y, whereas to say that \((\exists x)(\forall y)R\) is true means that there is a fixed object x such that R is a true relation between this object and any object y.

